% Contenido para: sections/2introduccion.tex

En el campo de la ingeniería eléctrica y la metrología, el osciloscopio se erige como una de las herramientas más fundamentales y versátiles para el análisis de circuitos. Este instrumento permite la visualización gráfica de señales eléctricas en el dominio del tiempo, ofreciendo una visión directa de parámetros cruciales como la amplitud, la frecuencia, el período, el desfase y la forma de onda. Comprender su operación no es solo un requisito académico, sino una competencia esencial para el diagnóstico, diseño y validación de cualquier sistema electrónico.

Esta experiencia de laboratorio se centra en la operación fundamental del osciloscopio digital. Se busca que el estudiante pase de un conocimiento nulo a una operación competente, dominando los controles básicos que permiten estabilizar y medir una señal. Se pondrá especial énfasis en los sistemas de disparo (trigger), los modos de acoplamiento y los ajustes de escala vertical (ganancia) y horizontal (base de tiempo), conceptos que constituyen la base para cualquier medición avanzada.

\subsection*{2.1 Objetivos}

\subsubsection*{Objetivos Generales}
\begin{itemize}
    \item Aprender a usar un osciloscopio y su funcionamiento.
    \item Conocer y aprender los conceptos de básicos de operación como son: base de tiempo, ganancia por canal, acoplamiento, trigger, etc.
    \item Conocer y aprender los principios de conexión a partir de las puntas o sondas de medida.
    \item Analizar los resultados y obtener conclusiones.
\end{itemize}

\subsubsection*{Objetivos Específicos}
Analizar el concepto y el comportamiento práctico de:
\begin{itemize}
    \item La base de tiempo y la ganancia por canal.
    \item El trigger y level.
    \item Los tipos de acoplamientos: CC, CA y GND.
    \item El modo normal e invertido.
    \item El modo Y-t y el modo X-Y.
\end{itemize}